% questions.tex —— 2012 年·全国大纲卷·文科 · 题目内容片段（无答案）
% 由 tools/generate_projects.py 依来源总账自 2012/2012_数学真题_全国大纲文_试题.md 生成，勿手改（改源 MD 后重跑）


\gkpart{数学}

\begin{gknotice}
\textbf{来源说明}：本题面依据固定源 PDF 的两页 A3 题面逐题转录，原卷印刷页码为 727–728。

\textbf{解析归属}：固定源不附答案或解析；配套答案与解析均由本库按本题面独立复算。

\textbf{图形说明}：第 19 题立体图由项目内 TikZ 矢量源重绘。下列 Markdown 使用逻辑 PNG 标记供后续生成器映射，仓库不保存 PNG。
\end{gknotice}

\gkpart{一、选择题}

\begin{question}
\gkti{1}已知集合 $A=\{x\mid x\text{ \text{是平行四边形}}\}$，$B=\{x\mid x\text{ \text{是矩形}}\}$，$C=\{x\mid x\text{ \text{是正方形}}\}$，$D=\{x\mid x\text{ \text{是菱形}}\}$，则\mbox{（　　）}

（A）$A\subseteq B$　（B）$C\subseteq B$　（C）$D\subseteq C$　（D）$A\subseteq D$

\end{question}
\begin{question}
\gkti{2}函数 $y=\sqrt{x+1}$（$x\geq-1$）的反函数为\mbox{（　　）}

（A）$y=x^2-1$（$x\geq0$）　（B）$y=x^2-1$（$x\geq1$）

（C）$y=x^2+1$（$x\geq0$）　（D）$y=x^2+1$（$x\geq1$）

\end{question}
\begin{question}
\gkti{3}若函数 $f(x)=\sin\dfrac{x+\varphi}{3}$（$\varphi\in[0,2\pi]$）是偶函数，则 $\varphi=$\mbox{（　　）}

（A）$\dfrac\pi2$　（B）$\dfrac{2\pi}3$　（C）$\dfrac{3\pi}2$　（D）$\dfrac{5\pi}3$

\end{question}
\begin{question}
\gkti{4}已知 $\alpha$ 为第二象限角，$\sin\alpha=\dfrac35$，则 $\sin2\alpha=$\mbox{（　　）}

（A）$-\dfrac{24}{25}$　（B）$-\dfrac{12}{25}$　（C）$\dfrac{12}{25}$　（D）$\dfrac{24}{25}$

\end{question}
\begin{question}
\gkti{5}椭圆的中心在原点，焦距为 $4$，一条准线为 $x=-4$，则该椭圆的方程为\mbox{（　　）}

（A）$\dfrac{x^2}{16}+\dfrac{y^2}{12}=1$　（B）$\dfrac{x^2}{12}+\dfrac{y^2}{8}=1$

（C）$\dfrac{x^2}{8}+\dfrac{y^2}{4}=1$　（D）$\dfrac{x^2}{12}+\dfrac{y^2}{4}=1$

\end{question}
\begin{question}
\gkti{6}已知数列 $\{a_n\}$ 的前 $n$ 项和为 $S_n$，$a_1=1$，$S_n=2a_{n+1}$，则 $S_n=$\mbox{（　　）}

（A）$2^{n-1}$　（B）$\left(\dfrac32\right)^{n-1}$　（C）$\left(\dfrac23\right)^{n-1}$　（D）$\dfrac1{2^{n-1}}$

\end{question}
\begin{question}
\gkti{7}$6$ 位选手依次演讲，其中选手甲不在第一个也不在最后一个演讲，则不同的演讲次序共有\mbox{（　　）}

（A）$240$ 种　（B）$360$ 种　（C）$480$ 种　（D）$720$ 种

\end{question}
\begin{question}
\gkti{8}已知正四棱柱 $ABCD-A_1B_1C_1D_1$ 中，$AB=2$，$CC_1=2\sqrt2$，$E$ 为 $CC_1$ 的中点，则直线 $AC_1$ 到平面 $BED$ 的距离为\mbox{（　　）}

（A）$2$　（B）$\sqrt3$　（C）$\sqrt2$　（D）$1$

\end{question}
\begin{question}
\gkti{9}$\triangle ABC$ 中，$AB$ 边的高为 $CD$，若 $\overrightarrow{CB}=\boldsymbol a$，$\overrightarrow{CA}=\boldsymbol b$，$\boldsymbol a\cdot\boldsymbol b=0$，$|\boldsymbol a|=1$，$|\boldsymbol b|=2$，则 $\overrightarrow{AD}=$\mbox{（　　）}

（A）$\dfrac13\boldsymbol a-\dfrac13\boldsymbol b$　（B）$\dfrac23\boldsymbol a-\dfrac23\boldsymbol b$　（C）$\dfrac35\boldsymbol a-\dfrac35\boldsymbol b$　（D）$\dfrac45\boldsymbol a-\dfrac45\boldsymbol b$

\end{question}
\begin{question}
\gkti{10}已知 $F_1,F_2$ 为双曲线 $C:x^2-y^2=2$ 的左、右焦点，点 $P$ 在 $C$ 上，$|PF_1|=2|PF_2|$，则 $\cos\angle F_1PF_2=$\mbox{（　　）}

（A）$\dfrac14$　（B）$\dfrac35$　（C）$\dfrac34$　（D）$\dfrac45$

\end{question}
\begin{question}
\gkti{11}已知 $x=\ln\pi$，$y=\log_5 2$，$z=\mathrm e^{-\frac12}$，则\mbox{（　　）}

（A）$x<y<z$　（B）$z<x<y$　（C）$z<y<x$　（D）$y<z<x$

\end{question}
\begin{question}
\gkti{12}正方形 $ABCD$ 的边长为 $1$，点 $E$ 在边 $AB$ 上，点 $F$ 在边 $BC$ 上，$AE=BF=\dfrac13$。动点 $P$ 从 $E$ 出发沿直线向 $F$ 运动，每当碰到正方形的边时反弹，反弹时反射角等于入射角，当点 $P$ 第一次碰到 $E$ 时，$P$ 与正方形的边碰撞的次数为\mbox{（　　）}

（A）$8$　（B）$6$　（C）$4$　（D）$3$

\end{question}
\gkpart{二、填空题}

\begin{question}
\gkti{13}$\left(x+\dfrac1{2x}\right)^8$ 的展开式中 $x^2$ 的系数为\underline{\hspace{2cm}}。

\end{question}
\begin{question}
\gkti{14}若 $x,y$ 满足约束条件

\nopagebreak[4]
\[
\begin{cases}
x-y+1\geq0,\\
x+y-3\leq0,\\
x+3y-3\geq0,
\end{cases}
\]

则 $z=3x-y$ 的最小值为\underline{\hspace{2cm}}。

\end{question}
\begin{question}
\gkti{15}当函数 $y=\sin x-\sqrt3\cos x$（$0\leq x<2\pi$）取得最大值时，$x=$\underline{\hspace{2cm}}。

\end{question}
\begin{question}
\gkti{16}已知正方体 $ABCD-A_1B_1C_1D_1$ 中，$E,F$ 分别为 $BB_1,CC_1$ 的中点，那么异面直线 $AE$ 与 $D_1F$ 所成角的余弦值为\underline{\hspace{2cm}}。

\end{question}
\gkpart{三、解答题}

\begin{question}
\gkti{17}$\triangle ABC$ 的内角 $A,B,C$ 成等差数列，其对边 $a,b,c$ 满足 $2b^2=3ac$，求 $A$。

\end{question}
\begin{question}
\gkti{18}已知数列 $\{a_n\}$ 中，$a_1=1$，前 $n$ 项和 $S_n=\dfrac{n+2}{3}a_n$。

（1）　求 $a_2,a_3$；

（2）　求 $\{a_n\}$ 的通项公式。

\Needspace{0.4\textheight}
\end{question}
\begin{question}
\gkti{19}如图，四棱锥 $P-ABCD$ 中，底面 $ABCD$ 为菱形，$PA\perp$ 底面 $ABCD$，$AC=2\sqrt2$，$PA=2$，$E$ 是 $PC$ 上的一点，$PE=2EC$。
\nopagebreak[4]
\begin{center}% 2012 全国大纲卷文科第 19 题：四棱锥 P-ABCD。
% 实线表示原图可见棱，虚线表示被遮挡的底面对角线或高线。
\begin{tikzpicture}[scale=.92, line cap=round, line join=round]
  \coordinate (P) at (2.35,4.05);
  \coordinate (A) at (2.35,1.45);
  \coordinate (B) at (0.35,0.80);
  \coordinate (C) at (1.20,0.05);
  \coordinate (D) at (4.40,0.75);
  \coordinate (E) at ($(P)!0.6667!(C)$);

  % 底面菱形与侧棱
  \draw[thick] (B)--(C)--(D);
  \draw[dashed] (D)--(A)--(B) (B)--(D);
  \draw[thick] (P)--(B) (P)--(D) (P)--(C);
  \draw[dashed] (P)--(A);

  % 点 E 在 PC 上；E 与 B、D 的连线给出截面 BED。
  \draw[thick] (E)--(B) (E)--(D);

  \fill (P) circle (1.15pt) node[above] {$P$};
  \fill (A) circle (1.05pt) node[right] {$A$};
  \fill (B) circle (1.05pt) node[left] {$B$};
  \fill (C) circle (1.05pt) node[below] {$C$};
  \fill (D) circle (1.05pt) node[right] {$D$};
  \fill (E) circle (1.05pt) node[left] {$E$};
\end{tikzpicture}
\end{center}
\nopagebreak[4]
（1）　证明：$PC\perp$ 平面 $BED$；

（2）　设二面角 $A-PB-C$ 为 $90^\circ$，求 $PD$ 与平面 $PBC$ 所成角的大小。

\end{question}
\begin{question}
\gkti{20}乒乓球比赛规则规定：一局比赛，双方比分在 $10$ 平前，一方连续发球 $2$ 次后，对方再连续发球 $2$ 次，依次轮换。每次发球，胜方得 $1$ 分，负方得 $0$ 分。设在甲、乙的比赛中，每次发球，发球方得 $1$ 分的概率为 $0.6$，各次发球的胜负结果相互独立。甲、乙的一局比赛中，甲先发球。

（1）　求开始第 $4$ 次发球时，甲、乙的比分为 $1$ 比 $2$ 的概率；

（2）　求开始第 $5$ 次发球时，甲得分领先的概率。

\end{question}
\begin{question}
\gkti{21}已知函数 $f(x)=\dfrac13x^3+x^2+ax$。

（1）　讨论 $f(x)$ 的单调性；

（2）　设 $f(x)$ 有两个极值点 $x_1,x_2$，若过两点 $(x_1,f(x_1))$，$(x_2,f(x_2))$ 的直线 $l$ 与 $x$ 轴的交点在曲线 $f(x)$ 上，求 $a$ 的值。

\end{question}
\begin{question}
\gkti{22}已知抛物线 $C:y=(x+1)^2$ 与圆 $M:(x-1)^2+\left(y-\dfrac12\right)^2=r^2$（$r>0$）有一个公共点 $A$，且在 $A$ 处两曲线的切线为同一直线 $l$。

（1）　求 $r$；

（2）　设 $m,n$ 是异于 $l$ 且与 $C$ 及 $M$ 都相切的两条直线，$m,n$ 的交点为 $D$，求 $D$ 到 $l$ 的距离。

\end{question}